Theorem 3 (Connection to the Shapley interaction index). Given a subset of input variables T ⊆ N, IShapley(T ) =∑
S⊆N\T
|S|!(|N|−|S|−|T| )!
(|N|−|T| +1)! ∆vT (xS) denotes the Shapley interaction index [26] of T . We have proved that the Shapley
interaction index can be represented as the weighted sum of causal effects involvingT , i.e.,IShapley(T ) =∑
S⊆N\T
1
|S|+1wS∪T .
In other words, the index IShapley(T ) can be explained as uniformly allocating causal effects wS′ s.t. S′ =S∪T to the
compositional variables ofS′, if we treat the coalition of variables inT as a single variable.
20
